\chapter{The conjecture}\label{ch:conjecture}

\section{Hodge classes}

Let $X$ be a smooth complex projective variety of dimension $d$. Its
cohomology with rational coefficients carries two structures that have, on
the face of it, nothing to do with each other. The first is the rational
structure itself, the lattice $H^{k}(X,\ZZ)$ modulo torsion inside
$H^{k}(X,\CC)$, which is a topological invariant. The second is the Hodge
decomposition
\begin{equation}\label{eq:hodgedec}
  H^{k}(X,\CC)=\bigoplus_{p+q=k}H^{p,q}(X),\qquad
  \overline{H^{p,q}(X)}=H^{q,p}(X),
\end{equation}
in which $H^{p,q}(X)$ is the space of classes represented by closed forms of
type $(p,q)$. The decomposition depends on the complex structure and moves
when the complex structure moves. The rational structure does not.

The Hodge classes are the classes where the two structures meet.

\begin{definition}
For $0\le p\le d$ the space of \emph{rational Hodge classes of degree $2p$} is
\[
  \Hdg^{p}(X)=H^{2p}(X,\QQ)\cap H^{p,p}(X),
\]
the intersection being taken inside $H^{2p}(X,\CC)$.
\end{definition}

Because the decomposition is preserved by cup product, in the sense that
$H^{p,q}\cup H^{p',q'}\subseteq H^{p+p',q+q'}$, the direct sum
$\Hdg^{*}(X)=\bigoplus_{p}\Hdg^{p}(X)$ is a graded $\QQ$-subalgebra of
$H^{\mathrm{even}}(X,\QQ)$. It is called the \emph{Hodge ring} of $X$.

There is no reason, a priori, for $\Hdg^{p}(X)$ to be anything but zero.
For a very general complex structure on a given manifold the space
$H^{p,p}$ is in general position with respect to the rational lattice, and
the intersection is as small as it can be. What forces it to be larger is
geometry.

\section{Classes of subvarieties}

Let $Z\subseteq X$ be a closed irreducible subvariety of codimension $p$.
Integration over the smooth part of $Z$ defines a linear form on closed
$(2d-2p)$-forms,
\[
  \alpha\longmapsto\int_{Z_{\mathrm{reg}}}\alpha ,
\]
which converges and vanishes on exact forms. This is a theorem of Lelong
for the convergence and of Stokes, applied on a resolution of $Z$, for the
vanishing. By Poincar\'e duality the form corresponds to a class
$[Z]\in H^{2p}(X,\QQ)$. In fact the class lies in the image of
$H^{2p}(X,\ZZ)$, since it is the Poincar\'e dual of the fundamental class
of $Z$ in Borel--Moore homology. Extending linearly to cycles, that is, to
formal combinations $\sum n_{i}Z_{i}$, and passing to rational
coefficients gives the \emph{cycle class map}
\[
  \cl^{p}\colon Z^{p}(X)_{\QQ}\longrightarrow H^{2p}(X,\QQ).
\]
It factors through rational equivalence, so it is defined on the Chow
group $\CH^{p}(X)_{\QQ}$.

\begin{proposition}\label{prop:classispp}
For every cycle $Z$ of codimension $p$ the class $\cl^{p}(Z)$ lies in
$\Hdg^{p}(X)$.
\end{proposition}

\begin{proof}
The class is rational by construction, so it suffices to show that it has
type $(p,p)$. Under Poincar\'e duality,
$H^{2p}(X,\CC)$ is dual to $H^{2d-2p}(X,\CC)$, and the summand
$H^{p,p}$ is the annihilator of every summand $H^{r,s}$ with $r+s=2d-2p$
except $H^{d-p,d-p}$. So we must show that $\int_{Z}\alpha=0$ for every
closed form $\alpha$ of type $(r,s)$ with $r+s=2d-2p$ and $r\ne d-p$. But
$Z_{\mathrm{reg}}$ is a complex manifold of dimension $d-p$, and the
restriction of a form of type $(r,s)$ to it is of type $(r,s)$ there. A form
of type $(r,s)$ on a complex manifold of dimension $d-p$ vanishes as soon
as $r>d-p$ or $s>d-p$. Since $r+s=2(d-p)$ and $r\ne d-p$, one of the two
inequalities holds.
\end{proof}

So subvarieties give Hodge classes. The conjecture says that they give all
of them.

\begin{conjecture}[Hodge]\label{conj:hodge}
Let $X$ be a smooth complex projective variety. For every $p$ the cycle
class map
\[
  \cl^{p}\colon\CH^{p}(X)_{\QQ}\longrightarrow\Hdg^{p}(X)
\]
is surjective.
\end{conjecture}

Hodge raised the question in his address to the International Congress of
1950, in a form with integral coefficients. The formulation above, with
rational coefficients, is the one that survived. The next
sections explain why each of its hypotheses is needed.

\section{Why rational coefficients}

One might ask whether every integral class of type $(p,p)$, that is, every
element of $H^{2p}(X,\ZZ)$ whose image in $H^{2p}(X,\CC)$ lies in
$H^{p,p}(X)$, is the class of an integral cycle. In degree two the answer is
yes; this is the content of Chapter~\ref{ch:lefschetz}. In higher degree it
is no.

Atiyah and Hirzebruch \cite{AH62} found torsion classes that are not
algebraic. Their argument is topological. The class of an algebraic cycle
is the pushforward of a fundamental class along a resolution, and such
classes are annihilated by certain odd-degree Steenrod operations. They
then exhibited smooth projective varieties with torsion classes, which are
automatically of type $(p,p)$, on which those operations do not vanish.
Totaro \cite{Tot97} later explained the obstruction through complex
cobordism, and found further examples.

Torsion classes are invisible rationally, so one might hope that the
integral statement holds modulo torsion. It does not. Koll\'ar
\cite{Kol92} showed that on a very general hypersurface $X\subset\PP^{4}$
of sufficiently large degree $m$, with $m$ divisible by a high power of a
prime, every curve has degree divisible by that power. Yet the generator of
$H^{4}(X,\ZZ)\cong\ZZ$, which is of type $(2,2)$ for dimension reasons, has
degree $1$. So the generator is not the class of a curve, although $m$
times it is, namely the class of the intersection of $X$ with a plane. Rationally
nothing is lost, and integrally the conjecture fails.

\section{Why projective}

The definitions make sense for every compact K\"ahler manifold, and the
Hodge decomposition holds there. One could ask whether every rational
$(p,p)$ class on a compact K\"ahler manifold is a combination of classes
of analytic subvarieties, or more generously of Chern classes of coherent
sheaves. For degree two, the theorem of Lefschetz holds in the K\"ahler
setting for classes that are Chern classes of line bundles. But a general
complex torus has no line bundles beyond the topologically trivial ones,
and no analytic subvarieties of positive dimension at all, while it may
still have Hodge classes in higher degree.

So the analytic version of the conjecture fails as soon as such a torus has
a nonzero Hodge class of degree four, and such tori exist.
Voisin \cite{Voi02b} went much further. She constructed complex tori, of Weil
type in the sense of Chapter~\ref{ch:weil}, carrying nonzero Hodge classes
of degree four, on which every coherent sheaf has vanishing Chern classes
in that degree. So the K\"ahler analogue fails even with Chern classes of
sheaves in place of subvarieties, and projectivity is used essentially.
We return to Voisin's tori in Section~\ref{sec:voisintori}, because they
are the same Weil classes this book is about, seen from the other side.

\section{What is known in general}

The conjecture is trivial in degrees $0$ and $2d$. In degree two it is the
theorem of Lefschetz (Chapter~\ref{ch:lefschetz}). The hard Lefschetz
theorem then gives it in degree $2d-2$ (Corollary~\ref{cor:degree2d-2}).
In particular it holds for every variety of dimension at most three. For
fourfolds the first open degree is four, and there it is open in general,
though known for many classes of fourfolds, among them uniruled ones
\cite{CM78}. Beyond that there is no general method. The known results are
proved family by family, and abelian varieties are the family about which
the most is known.

A useful reformulation is due to Mumford and Deligne. The rational
cohomology $H^{k}(X,\QQ)$ with its Hodge decomposition is a rational Hodge
structure of weight $k$, and every Hodge structure has a symmetry group,
the \emph{Mumford--Tate group}. It is the smallest algebraic subgroup of
$\GL(H^{k}(X,\QQ))$ defined over $\QQ$ whose real points contain the image
of the homomorphism that defines the Hodge decomposition
(Section~\ref{sec:mt}). Its variant the \emph{Hodge group} has the
property that the Hodge classes in tensor constructions on $H^{k}(X,\QQ)$
are exactly its invariants. For an abelian variety $A$ the whole
cohomology is the exterior algebra on $H^{1}(A,\QQ)$, and so the Hodge ring
of $A$ is the ring of invariants of a single reductive group acting on an
exterior algebra. That is why abelian varieties are tractable: the Hodge classes can
be computed. The difficulty, once they are computed, is to find the cycles.

\section{The plan}

The first half of the book makes the last two sentences precise.
Chapter~\ref{ch:toolkit} recalls the Hodge-theoretic tools,
Chapter~\ref{ch:lefschetz} proves the theorem in degree two, and
Chapter~\ref{ch:abelian} computes the cohomology of abelian varieties and
explains how the Hodge ring is governed by the Hodge group.

Chapter~\ref{ch:weil} introduces Weil's examples. Let $K$ be an imaginary
quadratic field acting on an abelian variety $A$ of dimension $2n$, in such
a way that the two eigenspaces of $K$ on the tangent space both have
dimension $n$. Then $H^{2n}(A,\QQ)$ contains a two-dimensional space of
Hodge classes, the Weil classes, which for a general such $A$ are not
products of divisor classes. We prove that for a general member of a family
of such varieties, the Hodge ring is spanned by powers of the polarisation
together with this plane. So the Hodge conjecture for the general member
is the algebraicity of one class.

Chapter~\ref{ch:known} surveys the cases in which this class is known to
be algebraic, and Chapter~\ref{ch:remaining} attempts the general case.
